\section{What changes in the draft}
\label{sec:draft}

Statements are grouped by whether the measurements support them. Nothing in
\code{sections/} was edited; this is a list, not a patch.

\subsection{Statements that survive}

\begin{itemize}
\item The arcminute FK amplitude. \code{insights.tex} line~193: ``at small
separation, FK reaches a few $\times10^{-6}$, about $0.5\%$ of Order-0 and
roughly twice the FF term''. Corrected: $+4.28\times10^{-6}$, $0.51\%$ of
Order-0, $2.1$ times FF. This is the number the draft leans on most and it
is unaffected, because at $0.5'$ the queried configuration really does sit
on the tabulated corner.

\item The selection rule and its mechanism. \code{insights.tex} lines
188--192, that FK concentrates in the modulus pair $\xi_\kappa$ and
$\xi_+$ and is suppressed in the other two. The corrected calculation
preserves the ordering: FK is $0.5$ to $1\%$ of Order-0 in the modulus pair
and $0.5$ to $0.8\%$ in the other two, having been $10^{-16}$ to $10^{-9}$
there before. The rule is a statement about which combinations FK enters
strongly, and that survives; the claim that it vanishes does not.

\item The channel algebra of \code{conclusion.tex} lines 50--58, that
$\xi_+$ is sourced by the sum
$\langle\Phi_{00}\Psi_+\Psi_+\rangle+\langle\Phi_{00}\Psi_\times\Psi_\times\rangle$
and $\xi_-$ by the difference. That is a statement about the vertex
structure and is untouched.

\item Figure~7, the driving-field $\zeta$ slices. It is built from
$\gamma$-resolved data and is correct as printed.

\item The parity null. $C_\ell^{EB}$ remains zero identically in the
corrected run, as it must by parity.

\item The Monte-Carlo validation, read as the draft states it: ``with the
input statistics held fixed''. It validates the diagram assembly. It is
structurally blind to this error, because it reads the same table.
\end{itemize}

\subsection{Statements that do not survive}

\begin{itemize}
\item \code{insights.tex} lines 194--196: FK ``overtakes the falling
Order-0 signal beyond $\sim2^\circ$''. There is no crossover. The corrected
ratio is flat at $0.5$ to $1\%$ across the whole range.

\item \code{insights.tex} lines 119--122, the caption of the draft's
two-point decomposition figure: ``overtakes Order-0 in absolute
amplitude beyond $\gamma\sim2^\circ$''. Same artefact.

\item \code{insights.tex} lines 196--199: in $\xi_-$ and
$\xi_{\kappa\gamma_t}$, ``FK drops to the numerical floor''. It does so at
$\gamma\to0$ only. At the plotted separations the corrected FK is $0.5$ to
$0.8\%$ of Order-0 in both.

\item \code{insights.tex} lines 202--204 and the caption of the
$C_\ell$ figure: ``the $\gamma$-flat FK term shows up as a low-$\ell$
excess''. The flatness is the interpolation weight. The excess, a factor 18
over Order-0 at $\ell=3$, becomes $-0.20$ with a resolved vertex; FK
instead appears as a steady $0.8\%$ correction over $39\le\ell\le800$.

\item \code{conclusion.tex} lines 34--36: the leakage ``lifts the
convergence and the parity-even shear at the few-percent level and
overtakes the linear signal in absolute amplitude beyond roughly
$2^\circ$''. The lift is a few tenths of a percent to one percent, not a
few percent, and there is no crossover.

\item \code{conclusion.tex} lines 60--62: ``The FK power therefore splits
equally between the two polarizations,
$\Delta C_\ell^{EE}=\Delta C_\ell^{BB}$''. The isotropy argument that
supports it is exact at zero separation, where $\Psi_+$ and $\Psi_\times$
are statistically equivalent. At finite separation the separation vector
picks out an axis and they are not, so the equality does not extend to the
plotted range. In the corrected run
$\max|\Delta C_\ell^{EE}| = 1.69\times10^{-6}$ against
$\max|\Delta C_\ell^{BB}| = 7.03\times10^{-7}$, that is $B/E\simeq0.42$ in
power.

\item \code{conclusion.tex} lines 133--136: ``the large-angle FK excess
sits at low multipoles ($\ell\lesssim10$)''. There is no large-angle
excess. The surrounding caveat about the relativistic Poisson source and
the observed-redshift correction remains worth making, but it is no longer
attached to a feature in the data.

\item \code{appendix.tex} lines 601--603, the caption of the validation
figure: ``FK its decay toward the sign change at $\gamma\simeq32^\circ$''.
That is where the four-neighbour set drops the $(1,1,1)$ node. It is an
interpolation switch described as physics.

\item \code{conclusion.tex} line 36: the leakage ``strengthens steadily
with source redshift''. This is not refuted here, it is untested. Two
separate effects bias it in exactly that direction, and both are addressed
in Section~\ref{sec:lambdamin}.
\end{itemize}

\subsection{Figures that must be replaced}

Six of the sixteen figures depend on the vertex.

\begin{table}[h]
\centering
\begin{tabular}{llp{6.6cm}}
\toprule
\# & figure & status \\
\midrule
2 & \code{analysis3\_NLO\_FFFK.pdf} & regenerated; FK no longer crosses
Order-0, and the two lower panels now carry a visible FK curve \\
3 & \code{analysis3\_cl\_full\_vs\_O0.pdf} & regenerated; the low-$\ell$
excess is gone \\
4 & \code{cl\_EB\_polarization.pdf} & regenerated; $EE$ and $BB$ no longer
coincide, $EB$ still zero \\
5 & \code{multiz\_kappa\_xi\_cl.pdf} & needs regeneration; its cached FK
slice comes from the same table. The corrected redshift dependence it
should show is measured in Section~\ref{sec:lambdamin} and plotted in
Figure~\ref{fig:redshift}, including the $\lammin$ bias on its lowest
slice \\
6 & \code{appendix\_mc\_workflow.pdf} & its FK curve moves with the
corrected table, but the agreement it demonstrates does not, because the
Monte-Carlo reads the same vertex through \code{driver\_stats.py} and both
sides move together. Regenerating it re-plots the corrected values; it does
not test the correction. Its caption separately needs the $32^\circ$
sentence removed, since that feature is the interpolation-weight zero \\
7 & \code{zeta\_driving\_field\_slices\_draft.pdf} & unchanged as printed,
being built from resolved data. Worth re-examining only for the cutoff
question of Section~\ref{sec:cutoff} \\
\bottomrule
\end{tabular}
\caption{The vertex-dependent figures. Figures 1 and 8 to 16 do not depend
on the vertex and are unaffected.}
\label{tab:figures}
\end{table}

Two cautions carried over from the figure manifest. The generator for
Figure~7 writes directly into the repository \code{figures/} directory with
a hardcoded path, so it overwrites the paper's file if run. Figure~5
depends on an external cache under \code{talk/}, so its FK slice must be
regenerated rather than reused. The regenerations here were produced by
\code{code/figures\_corrected/regenerate.py}, which imports each paper
generator, rebinds only its FK input and output paths, and calls its
\code{main()}. The plotting, styling and observable definitions are
therefore the paper's own, unmodified, and nothing was written into
\code{figures/}. The regenerated files are in
\code{driver\_field\_emulators/figures/corrected/}, named after the
generator's own output stem. The \code{\_permfix} versions are the ones to
use: they carry both the grid correction and the callable corrections of
Section~\ref{sec:callablefix}. The \code{\_corrected} versions carry the
grid correction only and are kept for comparison.
